\section{问题定义}

\subsection{训练图与候选关系}

令训练知识图谱为 $\mathcal{G}_{\mathrm{tr}}=(\mathcal{V},\mathcal{E}_{\mathrm{tr}})$，其中
$\mathcal{E}_{\mathrm{tr}}\subseteq \mathcal{V}\times\mathcal{R}_{\mathrm{tr}}\times\mathcal{V}$，
$\mathcal{R}_{\mathrm{tr}}$ 是训练关系词表。一个关系问答样本表示为
$q=(h,t,r^+)$，其中 $r^+$ 是目标关系，$h,t$ 是有序实体对。定义同一实体对在训练图中的已知关系集合为
\begin{equation}
    T_{\mathrm{tr}}(h,t)=\{r:(h,r,t)\in\mathcal{E}_{\mathrm{tr}}\}.
\end{equation}
候选池为
\begin{equation}
    P_q=\mathcal{R}_{\mathrm{tr}}\setminus\bigl(\{r^+\}\cup T_{\mathrm{tr}}(h,t)\bigr).
    \label{eq:pool}
\end{equation}
式~\ref{eq:pool} 只保证候选没有作为该实体对的\emph{已知训练关系}出现；它不声称候选事实在真实世界中为假。这一点是本文“反事实”术语的边界，也是开放世界知识图谱中避免过度解释的必要条件 \citep{ji2022survey,hogan2021kg}。

\subsection{两个评估目标}

本文区分两种答案选择规则。第一，模型在列出十个关系选项的题目上直接解码关系文本，评价 Exact Match（EM）和 OOV；它衡量模型是否保持原始 GRIP 的生成路径。第二，评估程序给定同一题目的 $\{r^+,r_1^-,\ldots,r_K^-\}$，用关系 continuation score $z_\theta(q,r)$ 逐项排序，评价 MRR、Hits@1/3/10 和 wrong-in-list。后者更直接测量候选 InfoNCE 是否改变关系边界，但不是模型直接解码所得的准确率。

给定冻结 bank $B(q)=\{r_1^-,\ldots,r_K^-\}$，候选损失写为
\begin{equation}
    \mathcal{L}_{\mathrm{cand}}(q)=-\log
    \frac{\exp(z_\theta(q,r^+)/\tau_c)}
    {\exp(z_\theta(q,r^+)/\tau_c)+\sum_{k=1}^K\exp(z_\theta(q,r_k^-)/\tau_c)}.
    \label{eq:infonce}
\end{equation}
其中 $\tau_c$ 是候选温度。训练阶段保留原始生成损失 $\mathcal{L}_{\mathrm{gen}}$，只加入
\begin{equation}
    \mathcal{L}=\mathcal{L}_{\mathrm{gen}}+\lambda_{\mathrm{cand}}\mathcal{L}_{\mathrm{cand}}.
    \label{eq:total}
\end{equation}
Stage 1 图上下文训练和 B1 generation-only 路径不变，使得匹配预算比较具有清晰的归因边界。
